% GENERATED FILE. Edit canonical lessons, then run npm run generate:books.
% canonical-source-hash: fnv1a64:cee91d45dcc3a448
% canonical-lessons: MR-C123-garaj, MR-C123-asa, MR-C123-svapna, MR-C123-magne, MR-C123-tharavne

\chapter{Need, Hope, Dream, To ask for, To decide}
\label{ch:mr-need-hope-dream-to-ask-for-to-decide}
\hlchaptermodality{123}

\begin{chapteropening}
\textbf{By the end of this chapter:} \emph{I can say a need, hope, a dream, to ask for and to decide.}

Complete the last lesson of chapter 123: I can say a need, hope, a dream, to ask for and to decide.
\end{chapteropening}

\section[\texorpdfstring{\mr{गरज} (garaj)}{गरज (garaj)}]{\mr{गरज} (garaj) — a need}

\label{lesson:MR-C123-garaj}



\begin{quote}
Before the new one: say the Marathi for wanted, then the Marathi for a wish.
\end{quote}

\subsection*{You'll want to know: \mr{गरज}}
\textbf{\mr{गरज}} — \emph{garaj} — \textquotedblleft{}a need\textquotedblright{}.

\begin{cousinweb}
\textbf{\mr{गरज}} (\emph{garaj}) is a loan from Arabic \emph{gharaḍ}, \textquotedblleft{}aim, purpose, need\textquotedblright{}, which reached Marathi through Persian. Its Arabic root is \emph{gh-r-ḍ} (Lane). Persian and Hindustani say \emph{gharaz}.
\end{cousinweb}

\begin{grammarlens}[title={What you've built}]
One more word for this chapter.
\end{grammarlens}

\subsection*{Your turn}
\begin{itemize}
\raggedright
  \item \emph{Say it:} \emph{garaj}
  \item \emph{Say it:} \emph{garaj}, once more
  \item \emph{Say it:} say \emph{garaj}
  \item \emph{Recall:} say the Marathi for wanted, then the Marathi for a wish, then say \emph{garaj} again
\end{itemize}

\begin{tcolorbox}[breakable,colback=teal!4,colframe=teal!35!black,title={Before you move on}]
What does \mr{गरज} mean? (\textquotedblleft{}A need\textquotedblright{}.) Say it once more.
\end{tcolorbox}
\section[\texorpdfstring{\mr{आशा} (āśā)}{आशा (āśā)}]{\mr{आशा} (āśā) — hope}

\label{lesson:MR-C123-asa}



\begin{quote}
Before the new one: say the Marathi for a wish, then the Marathi for a need.
\end{quote}

\subsection*{You'll want to know: \mr{आशा}}
\textbf{\mr{आशा}} — \emph{āśā} — \textquotedblleft{}hope\textquotedblright{}.

\begin{cousinweb}
\textbf{\mr{आशा}} (\emph{āśā}) is Sanskrit \emph{āśā}, \textquotedblleft{}wish, expectation\textquotedblright{}, taken over whole, a tatsama. CDIAL lists the inherited Marathi form as \emph{ās}, \textquotedblleft{}hope\textquotedblright{} (CDIAL 1456).
\end{cousinweb}

\begin{grammarlens}[title={What you've built}]
One more word for this chapter.
\end{grammarlens}

\subsection*{Your turn}
\begin{itemize}
\raggedright
  \item \emph{Say it:} \emph{āśā}
  \item \emph{Say it:} \emph{āśā}, once more
  \item \emph{Say it:} say \emph{āśā}
  \item \emph{Recall:} say the Marathi for a wish, then the Marathi for a need, then say \emph{āśā} again
\end{itemize}

\begin{tcolorbox}[breakable,colback=teal!4,colframe=teal!35!black,title={Before you move on}]
What does \mr{आशा} mean? (\textquotedblleft{}Hope\textquotedblright{}.) Say it once more.
\end{tcolorbox}
\section[\texorpdfstring{\mr{स्वप्न} (svapna)}{स्वप्न (svapna)}]{\mr{स्वप्न} (svapna) — a dream}

\label{lesson:MR-C123-svapna}



\begin{quote}
Before the new one: say the Marathi for a need, then the Marathi for hope.
\end{quote}

\subsection*{You'll want to know: \mr{स्वप्न}}
\textbf{\mr{स्वप्न}} — \emph{svapna} — \textquotedblleft{}a dream\textquotedblright{}.

\begin{cousinweb}
\textbf{\mr{स्वप्न}} (\emph{svapna}) is Sanskrit \emph{svapna}, \textquotedblleft{}sleep, dream\textquotedblright{}, taken over whole as a learned loan, a tatsama (CDIAL 13904). It is built on \emph{svap}, \textquotedblleft{}to sleep\textquotedblright{}.
\end{cousinweb}

\begin{grammarlens}[title={What you've built}]
One more word for this chapter.
\end{grammarlens}

\subsection*{Your turn}
\begin{itemize}
\raggedright
  \item \emph{Say it:} \emph{svapna}
  \item \emph{Say it:} \emph{svapna}, once more
  \item \emph{Say it:} say \emph{svapna}
  \item \emph{Recall:} say the Marathi for a need, then the Marathi for hope, then say \emph{svapna} again
\end{itemize}

\begin{tcolorbox}[breakable,colback=teal!4,colframe=teal!35!black,title={Before you move on}]
What does \mr{स्वप्न} mean? (\textquotedblleft{}A dream\textquotedblright{}.) Say it once more.
\end{tcolorbox}
\section[\texorpdfstring{\mr{मागणे} (māgṇe)}{मागणे (māgṇe)}]{\mr{मागणे} (māgṇe) — to ask for}

\label{lesson:MR-C123-magne}



\begin{quote}
Before the new one: say the Marathi for hope, then the Marathi for a dream.
\end{quote}

\subsection*{You'll want to know: \mr{मागणे}}
\textbf{\mr{मागणे}} — \emph{māgṇe} — \textquotedblleft{}to ask for\textquotedblright{}.

\begin{cousinweb}
\textbf{\mr{मागणे}} (\emph{māgṇe}) is a tadbhava, a word inherited from Sanskrit and reshaped by sound change. It comes from \emph{mārgati}, \textquotedblleft{}seeks, requests\textquotedblright{} (CDIAL 10074). It is built on \emph{mārga}, \textquotedblleft{}track, path\textquotedblright{}.
\end{cousinweb}

\begin{grammarlens}[title={What you've built}]
One more word for this chapter.
\end{grammarlens}

\subsection*{Your turn}
\begin{itemize}
\raggedright
  \item \emph{Say it:} \emph{māgṇe}
  \item \emph{Say it:} \emph{māgṇe}, once more
  \item \emph{Say it:} say \emph{māgṇe}
  \item \emph{Recall:} say the Marathi for hope, then the Marathi for a dream, then say \emph{māgṇe} again
\end{itemize}

\begin{tcolorbox}[breakable,colback=teal!4,colframe=teal!35!black,title={Before you move on}]
What does \mr{मागणे} mean? (\textquotedblleft{}To ask for\textquotedblright{}.) Say it once more.
\end{tcolorbox}
\section[\texorpdfstring{\mr{ठरवणे} (ṭharavṇe)}{ठरवणे (ṭharavṇe)}]{\mr{ठरवणे} (ṭharavṇe) — to decide}

\label{lesson:MR-C123-tharavne}



\begin{quote}
Before the new one: say the Marathi for a dream, then the Marathi for to ask for.
\end{quote}

\subsection*{You'll want to know: \mr{ठरवणे}}
\textbf{\mr{ठरवणे}} — \emph{ṭharavṇe} — \textquotedblleft{}to decide\textquotedblright{}.

\begin{cousinweb}
\textbf{\mr{ठरवणे}} (\emph{ṭharavṇe}) is the \textquotedblleft{}make\textquotedblright{} form of Marathi \emph{ṭharṇẽ}, \textquotedblleft{}to be settled\textquotedblright{}. CDIAL derives that verb from a reconstructed \emph{*stabhira}, \textquotedblleft{}fixed, firm\textquotedblright{}, rather than \emph{*stharati}, \textquotedblleft{}is firm\textquotedblright{}, and marks it as taken from Hindi (CDIAL 13680, 13742).
\end{cousinweb}

\begin{grammarlens}[title={What you've built}]
One more word for this chapter.
\end{grammarlens}

\subsection*{Your turn}
\begin{itemize}
\raggedright
  \item \emph{Say it:} \emph{ṭharavṇe}
  \item \emph{Say it:} \emph{ṭharavṇe}, once more
  \item \emph{Say it:} say \emph{ṭharavṇe}
  \item \emph{Recall:} say the Marathi for a dream, then the Marathi for to ask for, then say \emph{ṭharavṇe} again
\end{itemize}

\begin{tcolorbox}[breakable,colback=teal!4,colframe=teal!35!black,title={Before you move on}]
What does \mr{ठरवणे} mean? (\textquotedblleft{}To decide\textquotedblright{}.) Say it once more.
\end{tcolorbox}
